\section{Implementation Reference}
\label{app:implementation}
This appendix collects supporting transform and table-construction details
formerly repeated in the module manuals. The signal model and one-hop contract
remain in the main text;
Appendices~\ref{sec:legacy-transform} and~\ref{app:one-hop} preserve the
transform derivations and detailed scheduling argument. The user manuals
cover controls, display interpretation, and practical setting choices.

\subsection{Bit reversal and table reuse}
For a radix-2 complex transform of length $Q=2^p$, write an index as
$i=\sum_{b=0}^{p-1}i_b2^b$, where $i_b\in\{0,1\}$. Its bit-reversed index is
\begin{equation}
 R_p(i)=\sum_{b=0}^{p-1}i_b2^{p-1-b}.
 \label{eq:reverse-index}
\end{equation}
For example, $R_3(3)=6$: the binary index $011$ becomes $110$. Placing the
input at these addresses prepares the iterative decimation-in-time traversal
of Equation~\eqref{eq:butterfly}. Listing~\ref{lst:reverse-table} constructs
the table during initialization, not on each sample call.

\begin{lstlisting}[float=tbp,language=C++,caption={Reference construction of a bit-reversal table for $Q=2^p$.},label={lst:reverse-table}]
for (size_t i = 0; i < Q; ++i) {
    size_t x = i, reversed = 0;
    for (size_t bit = 0; bit < p; ++bit) {
        reversed = (reversed << 1) | (x & 1);
        x >>= 1;
    }
    reversal[i] = reversed;
}
\end{lstlisting}

The production real analyzer prepares a table for
$Q_{\max}=N_{\max}/2=2^{p_{\max}}$. A frame of real length $N$ needs a
complex transform of length $M=N/2=2^p$. For $i<M$, its table entry is
\begin{equation}
 R_p(i)=R_{p_{\max}}(i)\mathbin{\gg}(p_{\max}-p).
 \label{eq:reverse-reuse}
\end{equation}
The leading zero bits of the shorter index become trailing zeros after
reversal; the shift removes precisely those zeros. Each packing unit writes
its complex pair directly to that address. This fuses permutation with
packing and window application instead of requiring a separate array pass.

\subsection{Twiddle lookup and batch reference}
Let $T[j]=\exp(-\mathrm{i}2\pi j/N_{\max})$ denote the maximum-size twiddle
table. At a butterfly stage of span $s$ and pair index $a$, the required
factor is obtained by
\begin{equation}
 W_s^a=T[aN_{\max}/s].
 \label{eq:twiddle-reuse}
\end{equation}
All supported spans divide $N_{\max}$, so the index is integral. Real-spectrum
reconstruction similarly uses $W_N^k=T[kN_{\max}/N]$. Thus changing the
transform length does not require rebuilding trigonometric tables. Scalar
float or double precision is used for the corresponding scalar instantiation;
SIMD channels share the float table.

Listing~\ref{lst:batch-reference} gives the complete dependency order for a
windowed real transform. It is explanatory batch pseudocode, not a function
invoked in one production sample call. In the production scheduler the pair,
butterfly, and reconstruction loops are suspended between work quotas;
window-cache preparation is included in the pair units. The output and
smoothing stage follows as described in Section~\ref{sec:one-hop}.

\begin{lstlisting}[float=tbp,caption={Batch reference for the packed real transform. R is the bit-reversal table for M; butterfly updates use Equation~\eqref{eq:butterfly}, and reconstruction uses Equation~\eqref{eq:real} and its endpoint rules.},label={lst:batch-reference}]
M = N / 2
for n = 0, ..., M-1:
    z[R[n]] = u[2*n] + i*u[2*n+1]  // u is already windowed
for span = 2, 4, ..., M:
    for group = 0, span, ..., M-span:
        for pair = 0, ..., span/2-1:
            execute_butterfly(z, span, group, pair)
for k = 0, ..., M:
    X[k] = reconstruct_positive_bin(z, k)
return X[0], ..., X[M]
\end{lstlisting}

The omitted negative-frequency bins follow from
$X[N-k]=\overline{X[k]}$ for $1\le k<M$; neither module needs to materialize
them. The generic full-spectrum real-FFT class discussed earlier does so.
The positive-spectrum convention includes both DC and Nyquist.

\subsection{Window and smoothing conventions}
Windowing multiplies each selected input sample by its window weight before
packing. Tapered windows trade a wider main lobe for different sidelobe
behavior~\cite{harris1978}; denser FFT bins alone do not remove that tradeoff.
Production weights use the existing tabulated coherent-gain correction,
$\widetilde w[n]=w[n]/G_{\mathrm{tab}}$, with the finite-length caveat in
Section~\ref{sec:signal-model}. A frame selects its input endpoint and
settings once, so later parameter changes cannot mix window functions within
that frame.

For completeness, the inclusive fractional-octave interval used by the
prefix-sum smoother can be written explicitly. Let $\Delta f=f_s/N$,
$f_k=k\Delta f$, and $\beta>0$. Define
\begin{align}
 f_h&=\min(f_k2^{\beta/2},f_s/2),\\
 f_\ell&=\begin{cases}
 f_k2^{-\beta/2},&f_k2^{\beta/2}\le f_s/2,\\
 f_h2^{-\beta},&f_k2^{\beta/2}>f_s/2,
 \end{cases}\\
 \ell_k&=\floor{f_\ell/\Delta f},\qquad
 h_k=\min(N/2,\floor{f_h/\Delta f}).
 \label{eq:octave-bounds}
\end{align}
The upper-edge adjustment preserves the nominal frequency ratio before
rounding to bins. The mean includes both endpoints; this is why the prefix
formula in Appendix~\ref{app:one-hop} uses $P_{h_k+1}$ and divides by
$h_k-\ell_k+1$. With smoothing disabled, the magnitude passes through
without this range average. DC has the one-bin interval $[0,0]$.

